\section*{Thesis TODOs}

\begin{enumerate}
    \item Bias-table (historical: 1985-2014, mid: 2020-2050, late: 2070-2100, group into neu,wce,med european areas)
    \item Learning rate scheduler
    \item Look back at the ecmwf course on ML for weather
    \item Look into gamma-correction for precipitation
    \item Look into BGS from \autoref{bgs-paper}. Can we guide image generation in pixels where we are certain of output data. I.e. scale the condition to target size, fill blank with zero, then guide the sampler with the error between that and the model estimate.
    \item Look into EDMs (\href{\url{https://egusphere.copernicus.org/preprints/2026/egusphere-2026-1139/egusphere-2026-1139.pdf}}{here} and \href{\url{https://arxiv.org/pdf/2206.00364}}{here})
\end{enumerate}

\section*{Figure outs - ML}
\begin{enumerate}
    \item \st{What about non-square, non-$2^n$ sized image-targets? Interpolation? Special architecture? (Targets are 412x424, nearest $2^n$ is 512, $412/(2^3)=51.5$).} zero-padding is better than interpolation! Use 432 size, as that has 4 possible halfings. 
    \item \st{Normalization on weather data. Model does not learn scale, but since scale is important for climate how do we "rescale" (un-normalize) the data?} You have to unnormalize properly....
\end{enumerate}

\section*{Figure outs - Climate}
\begin{enumerate}
    \item Does it make sense to evaluate on number of heatwaves? This is obvious from the condition data, so the model would never outperform the quality of the input.
    \item How do we define heatwaves?
    \item Other weather phenomenon I should look at, e.g. cloudbursts?
    \item \st{Should I start without precipitation?} No, include precipitation.
\end{enumerate}

\section*{Decisions}
\begin{enumerate}
    \item Use realization 1 for training and validation, realization 2 and 3 for inference.
\end{enumerate}
